\documentclass[11pt]{amsart}
\usepackage[T1]{fontenc}
\usepackage{lmodern,amsmath,amssymb,amsthm,mathtools,booktabs,array,longtable}
\usepackage[margin=1in]{geometry}
\usepackage{microtype}
\usepackage[hidelinks]{hyperref}
\pdfinfoomitdate=1
\pdftrailerid{}
\pdfsuppressptexinfo=15
\newtheorem{theorem}{Theorem}[section]
\newtheorem{proposition}[theorem]{Proposition}
\newtheorem{lemma}[theorem]{Lemma}
\newtheorem{corollary}[theorem]{Corollary}
\theoremstyle{definition}
\newtheorem{definition}[theorem]{Definition}
\newtheorem{remark}[theorem]{Remark}
\newcommand{\E}{\mathbb E}
\newcommand{\Pp}{\mathbb P}
\newcommand{\R}{\mathbb R}
\newcommand{\N}{\mathbb N}
\newcommand{\cH}{\mathcal H}
\newcommand{\cF}{\mathcal F}
\newcommand{\TV}{\operatorname{TV}}
\newcommand{\law}{\operatorname{Law}}
\newcommand{\Var}{\operatorname{Var}}
\newcommand{\supp}{\operatorname{supp}}
\newcommand{\cp}{\mathrm{cp}}
\newcommand{\on}{\mathrm{on}}
\newcommand{\norm}[1]{\left\lVert#1\right\rVert}
\newcommand{\abs}[1]{\left|#1\right|}
\title[Acquired geometry and causal resource transfer]{General Theta Foundations I:\\Acquired Geometry and Causal Resource Transfer}
\author{Qian Qi}
\date{October 7, 2026}
\subjclass[2020]{Primary 62B05, 60G35; Secondary 94A29, 28A80, 93E11}
\keywords{causal experiment, predictive quotient, actual acquisition, recursive quantization, intermittent mixing, singular measure}
\begin{document}
\begin{abstract}
An experiment supplies a law of acquired predictions, not a prescribed geometric parameter space. We prove a causal approximation theorem in which a finite recursive encoder is constructed from static inward covers and two inequalities for the raw report kernel. A Lyapunov envelope controls the locations visited by the approximate state; a two-point moment inequality controls its separation from the true predictive state. Their combination gives a uniform online error bound without a supplied recursive quantizer, a common invariant reference, or domination of intermediate strata by terminal masses. Actual subprobability small-ball estimates yield matching checkpoint lower bounds. A conditional-projection identity retains finite acquisition when the acquired law is atomic. We verify the theorem for an intermittently mixing hidden-state experiment with arbitrarily long informative noncontracting runs, and for a nonhomogeneous singular experiment transporting unrevealed tails through expanding changes of mode. The former has a matched label--output-precision--calibration law; the latter has an explicit acquired-depth law, including creation and disappearance of mode masses. Parameter-independent causal morphisms compose the same task risk with separately accounted simulation, controller, clock, workspace, and program resources. The results concern fixed exploration and deterministic horizons; a complete optimal region for all resource coordinates is not asserted.
\end{abstract}
\maketitle
\section{The experimental problem}\label{sec:problem}
A preparation and a sequence of causal interventions determine which predictions have actually been acquired. Two different obstructions then arise. A formally reachable set may carry little acquired mass, and a good terminal codebook may have no equally good recursive implementation. The first obstruction concerns probability; the second concerns the loss of access to discarded reports. A foundations theorem must connect them without replacing either by a dimension count.

The central result below is Theorem~\ref{thm:inward}. It starts with a report kernel and a predictive update, constructs the finite-label machine from a static covering, and controls that machine under the \emph{true} report law. This last qualification matters: a drift inequality under the approximate filter's own report distribution does not control an approximation driven by physical data. The construction combines inward quantization with two raw moment inequalities. It does not assume an error recurrence for an already constructed encoder. It also avoids comparison of all intermediate occupation masses with a terminal reference law.

The geometry is supplied by the actual terminal score law. Theorem~\ref{thm:inward} gives an upper bound on a general metric state domain. Proposition~\ref{prop:submass} turns a small-ball estimate for only a positive acquired submeasure into a matching lower bound, allowing the rest of the acquired law to be singular. Proposition~\ref{prop:oracle} treats the opposite difficulty: a finite acquisition protocol can make the entire predictive law atomic even though the unresolved physical quantity has a continuous or singular law. Its exact decomposition prevents substituting an ideal law for the actual acquisition.

Two verifications carry the general theorem beyond uniform one-step contraction. Theorem~\ref{thm:filter} treats a hidden-state experiment with observed intermittent mixing. Informative identity-transition observations are isometries in log odds and can occur in arbitrarily long runs. A finite inward compander on the unbounded log-odds domain nevertheless attains the optimal label scale without a logarithmic loss. Theorem~\ref{thm:singular} treats finite acquisition in a nonhomogeneous Cantor experiment. Both mode changes transport an old unrevealed tail, deletion expands errors, and mode probabilities may vary with time or vanish. No invariant preparation comparison enters the proof. The finite acquired law is computed explicitly.

The two results instantiate the same kernel theorem; neither is a premise of it. The complementary allocation-weighted, finite-chart, reset-preserved, and continuation-information results are retained in the accompanying mathematical supplement with their original hypotheses. The present route is not a claim that one stability condition subsumes every experiment. It supplies a different constructive mechanism: direct raw-kernel control of a recursively retained state, with acquired geometry supplying the converse.

\subsection{Raw experiments and legal information}
An experiment consists of standard Borel spaces $X_t,A_t,Y_{t+1}$, a probability preparation $\mu$ on $X_0$, and jointly measurable nonnegative normalized kernels
\begin{equation}\label{eq:raw}
 K_t(dx',dy,dc\mid x,a).
\end{equation}
Here $c$ records the declared nonnegative resource increments and elapsed physical time. Positivity means nonnegativity and normalization, not a uniform lower bound on report probabilities. Dirac kernels and failed acquisitions are allowed. A legal strategy is a kernel $\beta_t(da\mid h_t)$ using only the visible history, known calibration, and independent randomization. The history includes actions, reports, failures, and every published cost. A clock or remaining budget that changes legal continuations is part of the visible interface.

The Bayesian assertions fix the preparation and a legal exploration strategy independently of the encoder. An unknown parameter can be included in the hidden preparation with a specified prior; this does not turn a prior-specific quotient into a parameter-uniform sufficient statistic. The comparison theorem in Section~\ref{sec:resources} has its own uniform-in-parameter quantifiers. We work at a fixed deterministic horizon. Independent exact holds are allowed as specified below. Stopping rules are legal objects of the experiment, but the main risk theorem is not silently extended to arbitrary stopping times.

The successive kernel integrals define the actual path law. Write $\cH_t=\sigma(h_t)$ for its visible filtration. An online encoder retains a label in an alphabet of cardinality at most $M$, updates it using the current report and declared current action, and cannot reread earlier reports. A checkpoint encoder may inspect $h_T$ once and then retain at most $M$ labels. Independent randomization is allowed in both classes. Persistent labels do not include a free history tape, mode register, or real-valued posterior register. Required controller and internal clock states are counted separately in a serial implementation and are not silently removed from its total state.

\subsection{Executable tests and a realized quotient}
An executable future test is a legal finite continuation followed by a bounded measurable function of its actual reports. A hidden-state function is a test only when the experiment includes a measurement producing the required readout. Two histories are predictively equivalent when they have the same continuation interface and every such conditional test expectation agrees. A countable determining family $f_1,f_2,\ldots$ yields the evaluation map
\begin{equation}\label{eq:eval}
 h\longmapsto \big(\E[f_n\mid h]\big)_{n\ge1}.
\end{equation}
A realization requires measurability of its image and a valid update, not just a formal equivalence relation. The following sufficient certificate makes the relevant version issue explicit.

\begin{proposition}[Compact continuous-instrument realization]\label{prop:quotient}
Let $\mathcal B$ be a compact metrizable domain of beliefs, and let $\mathcal V\subset C(\mathcal B;\R)$ be the separable closed linear span of the continuous executable-test evaluations and declared continuation-interface coordinates, including constants. Choose a countable uniformly dense family in $\mathcal V$ for the evaluation map. Suppose, for each declared action $a$, there is a full-support reference measure $m_a$ on a Polish report domain and a continuous nonnegative density $g_a(\pi,y)$. On the declared domain $D_a=\{(\pi,y):g_a(\pi,y)>0\}$ let the normalized belief update $B_a(\pi,y)$ be continuous. Assume that, for every $f\in\mathcal V$ and every declared $y$, the function
\[
 \pi\longmapsto g_a(\pi,y) f(B_a(\pi,y))
\]
extends continuously by its instrument value to $\mathcal B$ and belongs to $\mathcal V$. Suppose these continuous evaluations determine all executable tests, including their allowed continuations. Then the image $S$ of the evaluation map on $\mathcal B$ is compact metrizable, predictive equivalence is equality in $S$, and the report law and update descend to $S$. The update is pointwise well defined on the positive-density declared domain. If a measurable statistic $r:H\to Z$, with $Z$ standard Borel, measurably determines every evaluation, the realized quotient is a measurable function of $r$ on its image (and after any fixed measurable extension off that image).
\end{proposition}
\begin{proof}
Normalize the countable determining evaluations to be bounded by one. Their product map $E:\mathcal B\to[-1,1]^{\N}$ is continuous, so $S=E(\mathcal B)$ is compact metrizable and Borel. Its fibers agree on the dense determining set, hence on $\mathcal V$ by uniform approximation, and hence on the declared executable tests. Taking $f=1$ in the instrument closure shows that equal fibers have equal $g_a(\pi,y)$ for every declared report. For positive common density, equality of the density-weighted evaluations followed by division by that \emph{positive} number shows that their updated beliefs have equal $E$-images. Thus both the report kernel and normalized update descend. Continuity of a map constant on fibers of a compact-to-Hausdorff quotient supplies continuity of the descended instrument evaluations. On each positive-density neighborhood the same quotient argument applies to the update. Nothing is asserted by normalization at density zero; an extra prescribed null-report version would require an extra closure certificate.

For the last assertion, for each $n$ let $u_n:Z\to\R$ be a measurable factor for the $n$th evaluation. The product $u=(u_n)_n$ is measurable. Since $S$ is Borel, replace $u(z)$ by a fixed $s_0\in S$ whenever $u(z)\notin S$. This is a measurable $S$-valued extension and agrees with $E$ on all histories. Under completed, almost-sure conventions the countably many equalities are first imposed on one common full-measure history set. No factorization through an unspecified arbitrary measurable codomain is used.
\end{proof}

The raw realizations below also verify their quotients directly: a finite-dimensional posterior, or a mode and cylinder posterior, determines every legal continuation and is separated by a terminal test. The compact certificate is sufficient, not necessary; the stable metric used for quantization need not induce the compact quotient topology. In particular, log odds has an unbounded metric even though a finite-state posterior lies in a compact simplex.

\subsection{A common squared task}
Fix a scalar terminal report $Y\in[0,1]$, or a finite collection of separately executable scalar probes with prescribed positive weights. Expectations of the weighted sum of their squared losses define the vector task; no joint counterfactual observation is postulated. Use the corresponding Euclidean norm. Set
\begin{equation}\label{eq:score}
 P_T=\E[Y\mid\cH_T],\qquad \nu_T=\law(P_T),\qquad
 B_T=\E\norm{Y-P_T}^2.
\end{equation}
All statements below are componentwise for this interpretation. For a probability law $\nu$ on the task box define
\begin{equation}\label{eq:q}
 q_M(\nu)=\inf_{\#C\le M}\int \operatorname{dist}(p,C)^2\,\nu(dp).
\end{equation}
The infimum permits arbitrary centers; projection to the task box is harmless. The symbols $R_\cp(T,M)$ and $R_\on(T,M)$ denote the infimal risks in the two legal classes. The law $\nu_T$ always comes from the actual preparation and exploration, not an invariant reference chosen for analytic convenience.

\section{The causal acquired-geometry theorem}\label{sec:transfer}
Let $S$ be a standard Borel predictive-state domain equipped with a separable metric $d$ generating its Borel structure. The physical state $S_t$ here is the \emph{predictive} state derived from the experiment, not necessarily its hidden microscopic state. At an active step, its report kernel and update are
\[
 J_t(dy\mid x),\qquad F_t(x,y)\in S.
\]
Both are Borel. The update is declared on all true and representative states used below, including the report versions needed to evaluate a representative on a physical report. This is a genuine hypothesis: almost-sure versions separately chosen for each representative do not suffice. The state and the strategy interface determine the true report kernel; when a fixed controller is needed it is included in the declared interface and separately charged.

Let $w:S\to[1,\infty)$ be a Borel envelope. Static covering data at resolution $r_M>0$ consist of at most $M$ representatives $c_1,\ldots,c_m$ and a Borel cover $A_1,\ldots,A_m$ of $S$ satisfying
\begin{equation}\label{eq:inward}
 x\in A_i\quad\Longrightarrow\quad
 d(x,c_i)\le r_M w(x),\qquad w(c_i)\le w(x).
\end{equation}
No recursive coding rule is included in this assumption. In an exact Borel machine membership in the finite cover is a declared operation. A finite-bit implementation additionally requires certified membership or enclosure routines; its error and workspace are not inferred from Borel measurability alone.

Let $\cF_{t-1}$ contain the entire observed action--report past, the encoder's retained label, and its independent randomization already used. It does not contain unobserved hidden states or latent apparatus coins. Independent public randomization is allowed but is not an information source about the preparation. Let $D_t\in\{0,1\}$ be the active-step indicator. For some deterministic $\vartheta_t\in[0,1]$, require
\begin{equation}\label{eq:conditional}
 \Pp(D_t=1\mid\cF_{t-1})=\vartheta_t,\qquad
 \law(Y_t\mid\cF_{t-1},D_t=1)=J_t(\cdot\mid S_{t-1}).
\end{equation}
On $D_t=0$ both the predictive state and retained label are copied exactly. The hold report and its cost are still recorded. Any active randomization used to choose $F_t$ is included in $Y_t$. The first identity, being constant conditionally on the entire past, is the precise independence needed for the hold mixture. No independence between $Y_t$ and the old approximation error is assumed.

The raw moment conditions, uniform over the indicated deterministic times and all $x,z\in S$, are
\begin{align}
 \int d(F_t(x,y),F_t(z,y))^2 J_t(dy\mid x)
   &\le \kappa d(x,z)^2,\qquad 0\le\kappa<1,\label{eq:pair}\\
 \int w(F_t(z,y))^2 J_t(dy\mid x)
   &\le a w(z)^2+b,\qquad 0\le a<1,
       \quad 0\le b<\infty.\label{eq:envelope}
\end{align}
Crucially, the integration in \eqref{eq:envelope} is under the true state $x$, even though the envelope is evaluated at the candidate update from $z$. No stationary law or terminal occupation domination occurs in these conditions. They permit occasional expansion, but are sufficient conditions rather than a characterization of all finite-memory predictable experiments.

\begin{theorem}[Raw-kernel construction and acquired-risk transfer]\label{thm:inward}
Fix a prepared causal experiment, a deterministic horizon $T$, and exploration independent of the encoder. Suppose its realized predictive state satisfies \eqref{eq:conditional}--\eqref{eq:envelope}, and static data \eqref{eq:inward} are given for each allowed label budget $M$. Suppose a legal seed routine retains a representative $\widehat S_0$ in this codebook with
\[
 \E w(\widehat S_0)^2\le H_0,
 \qquad \norm{d(S_0,\widehat S_0)}_{L^2}\le e_0.
\]
All seed calls and states are charged. A known deterministic prepared state is an admissible seed. Suppose the task score is $P_T=f_T(S_T)$ and $f_T$ is $L$-Lipschitz for the task norm. Define
\begin{equation}\label{eq:radius}
 H=\max\{H_0,b/(1-a)\},\qquad
 E_M=\max\left\{e_0,\frac{r_M\sqrt H}{1-\sqrt\kappa}\right\}.
\end{equation}
Then the static data determine a Borel recursive encoder with at most $M$ retained labels, using only its preceding label and the current active report, such that
\begin{equation}\label{eq:main}
 B_T+q_M(\nu_T)=R_\cp(T,M)
 \ \le\ R_\on(T,M)
 \ \le\ B_T+L^2E_M^2.
\end{equation}
The constants are independent of $T$ and the sequence $(\vartheta_t)$ whenever the displayed hypotheses have uniform constants. They contain no reciprocal terminal stratum mass. If an actual submeasure $\underline\nu_T\le\nu_T$ of mass $m_T>0$ satisfies, for a declared $s_M>0$,
\begin{equation}\label{eq:sb-main}
 \sup_{u}\underline\nu_T(B(u,s_M))\le \frac{m_T}{2M},
\end{equation}
then $R_\cp(T,M)-B_T\ge m_T s_M^2/2$. In particular, if $e_0\le C_0r_M$, $s_M\ge c_0r_M$, and $m_T$ is bounded below, checkpoint and online excess risks have the same order $r_M^2$ on this class.
\end{theorem}

The theorem separates three distinct inputs: a static global cover, local inequalities for the raw report kernel, and a converse from actual acquired mass. None is an assumption about the optimal online risk. A complicated model may still make any of these inputs difficult to verify. Its advantage over a supplied-local-machine hypothesis is logical and constructive: the finite recursive machine and its error recurrence are conclusions. Its advantage over an occupation-reference argument is that transient changes in support do not create uncharged error terms, provided the uniform two-point and envelope bounds hold.

\begin{lemma}[Conditional projection and checkpoint quantization]\label{lem:projection}
For the task in \eqref{eq:score}, every forecast $a$ based on the history and independent algorithm randomization satisfies
\begin{equation}\label{eq:projection}
 \E\norm{Y-a}^2=B_T+\E\norm{P_T-a}^2.
\end{equation}
Consequently $R_\cp(T,M)=B_T+q_M(\nu_T)$, including independent shared-seed and randomized encoders or decoders.
\end{lemma}
\begin{proof}
Condition on the history and algorithm randomization. The conditional mean of $Y$ is still $P_T$, so the cross term vanishes. Fix a value of an independent shared seed. A randomized decoder can be replaced by its conditional mean given the retained label, leaving at most $M$ centers and decreasing squared loss. For any such center list, randomized label selection cannot improve on the closest center. Choose the first closest center to obtain a Borel selector. Thus every seed-conditional distortion is at least $q_M(\nu_T)$; independence of the seed leaves the law $\nu_T$ unchanged. Conversely any finite list and its nearest-center checkpoint selector approach the infimum in \eqref{eq:q}. Integration over the seed proves the assertion. The same argument applies separately to each executable probe.
\end{proof}

\begin{lemma}[Inward envelope and true-law error recursion]\label{lem:recursion}
Choose the first index $i$ for which $x\in A_i$ and put $Q_M(x)=c_i$. On active steps define
\begin{equation}\label{eq:machine}
 \widetilde S_t=F_t(\widehat S_{t-1},Y_t),\qquad
 \widehat S_t=Q_M(\widetilde S_t);
\end{equation}
on holds copy $\widehat S_{t-1}$. Then $\E w(\widehat S_t)^2\le H$ for every $t$, and with $e_t=\norm{d(S_t,\widehat S_t)}_{L^2}$,
\begin{equation}\label{eq:energy}
 e_t^2\le(1-\vartheta_t)e_{t-1}^2+
 \vartheta_t\big(\sqrt\kappa e_{t-1}+r_M\sqrt H\big)^2.
\end{equation}
In particular $e_t\le E_M$.
\end{lemma}
\begin{proof}
The least-index selection in a finite Borel cover is Borel and obeys both inequalities in \eqref{eq:inward}. Given $\cF_{t-1}$, both the true and approximate states are known, so \eqref{eq:conditional} and \eqref{eq:envelope} imply
\[
 \E[w(\widehat S_t)^2\mid\cF_{t-1},D_t=1]
 \le a w(\widehat S_{t-1})^2+b.
\]
Induction and $aH+b\le H$ yield the envelope bound, after mixing with the exact hold branch. The same bound holds for $\E[w(\widetilde S_t)^2\mid D_t=1]$, since $D_t$ is independent of the previous complete past. If $\vartheta_t=0$ only the hold branch is present and no conditional-on-active assertion is needed.

On the active branch, the triangle inequality and \eqref{eq:inward} give
\[
 d(S_t,\widehat S_t)
 \le d(F_t(S_{t-1},Y_t),F_t(\widehat S_{t-1},Y_t))
       +r_Mw(\widetilde S_t).
\]
Use \eqref{eq:pair} conditionally on the complete past and then Minkowski in $L^2$ on the active branch. The first term has norm at most $\sqrt\kappa e_{t-1}$ and the second at most $r_M\sqrt H$. On holds the old error is unchanged. Mixing the \emph{squares} of these branch norms proves \eqref{eq:energy}. Finally, if $e_{t-1}\le E_M$, the active root is at most $E_M$ by \eqref{eq:radius}; both branch squares are then at most $E_M^2$. Induction starts at the stated seed. This also proves the bound when there are arbitrarily long intervals with no active steps.
\end{proof}

\begin{proof}[Proof of Theorem~\ref{thm:inward}]
The rule \eqref{eq:machine} retains only its representative index, and hence is a legal $M$-label recursion. Its readout is $f_T(\widehat S_T)$, projected to the task box if necessary. Lemma~\ref{lem:recursion} and Lipschitz continuity give root-mean-square task error at most $LE_M$. Lemma~\ref{lem:projection} gives the two equal-baseline bounds in \eqref{eq:main}; every online output is also a legal checkpoint output, which gives the middle comparison. For any list of at most $M$ centers, the union of the open $s_M$-balls has $\underline\nu_T$-mass at most $m_T/2$. Outside the union the distance to the list is at least $s_M$, including boundary points. Its distortion is therefore at least $m_Ts_M^2/2$. This proves the last assertion, also for randomized encoders by Lemma~\ref{lem:projection}.
\end{proof}

\begin{corollary}[Controlled numerical perturbations]\label{cor:perturb}
In Theorem~\ref{thm:inward}, replace the active candidate $F_t(z,y)$ by a declared Borel numerical value $G_t(z,y)$ satisfying
\begin{equation}\label{eq:perturb}
 d(G_t(z,y),F_t(z,y))\le u,
 \qquad w(G_t(z,y))\le w(F_t(z,y))
\end{equation}
for all declared inputs. Allow an additional root-mean-square readout error $v$. Then
\[
 R_\on(T,M)\le B_T+(L E_{M,u}+v)^2,
 \qquad
 E_{M,u}=\max\left\{e_0,
        \frac{r_M\sqrt H+u}{1-\sqrt\kappa}\right\}.
\]
A deterministic bound $u=u_{\rm cal}+u_{\rm num}+u_{\rm input}$ may combine calibration, arithmetic, and finite-input enclosure errors when \eqref{eq:perturb} has been certified for their common implementation.
\end{corollary}
\begin{proof}
The inward envelope inequality in \eqref{eq:perturb} preserves \eqref{eq:envelope} for the numerical candidate and therefore preserves the proof of the envelope bound. The active error triangle gains the deterministic term $u$ before Minkowski is applied. Replace $r_M\sqrt H$ by $r_M\sqrt H+u$ in \eqref{eq:energy} and its invariant radius. The additional readout perturbation adds in the same $L^2$ task norm. No assertion about arbitrary, envelope-increasing rounding is used.
\end{proof}

\section{Acquisition, geometric converses, and output cuts}\label{sec:geometry}
The geometric quantity in \eqref{eq:main} is the quantization error of the actual score law. The following elementary consequences specify which kinds of acquired evidence suffice for a lower bound and how finite acquisition changes that evidence.

\begin{proposition}[A partial acquired-mass witness]\label{prop:submass}
Let $\underline\nu\le\nu$ have mass $m>0$. Suppose
\[
 \sup_u\underline\nu(B(u,r))\le A r^d
 \quad (0<r\le r_0),\qquad d>0.
\]
For any $M$ with $(m/(2AM))^{1/d}\le r_0$,
\[
 q_M(\nu)\ge \frac m2\left(\frac{m}{2AM}\right)^{2/d}.
\]
The submeasure may be the joint law of the score and a positive-probability event. No regularity is required of $\nu-\underline\nu$.
\end{proposition}
\begin{proof}
Apply \eqref{eq:sb-main} with $s_M=(m/(2AM))^{1/d}$. The union argument in the proof of Theorem~\ref{thm:inward} uses only domination $\underline\nu\le\nu$, not a decomposition into observable or separated strata.
\end{proof}

\begin{proposition}[Actual acquisition and a common oracle]\label{prop:oracle}
Let $\mathcal H\subset\mathcal G$ be sigma fields under a fixed prepared experiment, let $Y$ be the bounded scored report, and set
\[
 X=\E[Y\mid\mathcal G],\quad P=\E[X\mid\mathcal H],
 \quad b_\circ=\E\norm{Y-X}^2,\quad
 A=\E\norm{X-P}^2,\quad \nu=\law(P).
\]
An encoder may use $\mathcal H$ but not the additional oracle information in $\mathcal G$. Then
\begin{equation}\label{eq:oracle}
 R_\cp(M)-b_\circ=A+q_M(\nu),\qquad
 A+q_M(\nu)\ge q_M(\law(X)).
\end{equation}
If a legal online construction has $R_\on(M)-b_\circ\le A+C r_M^2$ and $q_M(\law(X))\ge c r_M^2$, then
\begin{equation}\label{eq:oracle-match}
 A+q_M(\nu)\asymp R_\on(M)-b_\circ\asymp A+r_M^2,
\end{equation}
with constants depending only on $c,C$. The statement allows $\nu$ to be finite atomic.
\end{proposition}
\begin{proof}
Successive conditional projection gives, for each forecast $a$ based on $\mathcal H$ and independent randomization,
\[
 \E\norm{Y-a}^2=b_\circ+A+\E\norm{P-a}^2
                  =b_\circ+\E\norm{X-a}^2.
\]
Minimize the first identity by Lemma~\ref{lem:projection}. After fixing any independent shared seed and replacing randomized readouts by their label-conditional means, the second distortion uses at most $M$ centers for $X$ and is bounded below by its unrestricted quantization error. Hence the checkpoint excess is at least both $A$ and $cr_M^2$, and therefore at least $\tfrac12\min(1,c)(A+r_M^2)$. The stated online upper and the checkpoint--online order prove \eqref{eq:oracle-match}. The ideal law supplies a lower witness, not a replacement for the actual law in the exact identity.
\end{proof}

\begin{proposition}[Persistent-label and output-alphabet cuts]\label{prop:output}
Let $G$ be a fixed finite subset of the task box, independent of the acquired history. Require the final physical forecast to lie in $G$, while the encoder retains at most $M$ labels. Then the nominal checkpoint excess over $B_T$ is
\begin{equation}\label{eq:qgrid}
 q_{M,G}(\nu_T)
  :=\min_{C\subset G,\ \#C\le M}\int\operatorname{dist}(p,C)^2\nu_T(dp),
\end{equation}
and in particular it is at least $q_{\min(M,\#G)}(\nu_T)$. If a nominal online readout has root-mean-square score error $e$ and every point of the task box lies within $h$ of $G$, rounding gives risk at most $B_T+(e+h)^2$.
\end{proposition}
\begin{proof}
For a fixed encoder, the decoder's expected loss is affine in its probability distribution over $G$, label by label; an extreme point minimizes it. It is therefore enough to consider deterministic decoders with at most $\min(M,\#G)$ distinct outputs. Nearest-center checkpoint selection proves \eqref{eq:qgrid}. Conditioning on an independent shared seed gives the same lower bound. Rounding the online readout changes it by at most $h$, and Minkowski followed by conditional projection gives the upper. This is a restriction on the output alphabet, not a theorem about every use of numerical precision.
\end{proof}

\begin{lemma}[A common calibration ambiguity]\label{lem:calibration}
Let a bounded latent nominal score $X$ satisfy $\E X=1/2$ and $X\pm\delta\in[0,1]$. The acquisition history and all its legal processing have laws independent of $\theta\in\{-\delta,\delta\}$, and the terminal report has conditional mean $X+\theta$. Let
\[
 b_\delta=\E[X(1-X)]-\delta^2,
 \quad P=\E[X\mid\mathcal H],\quad A=\E(X-P)^2.
\]
For every forecast $a$ independent of the unknown sign,
\begin{equation}\label{eq:calibration}
 \sup_{\theta=\pm\delta}\E_\theta(Y-a)^2-b_\delta
 =A+\E(P-a)^2+\delta^2+2\delta\abs{\E(P-a)}.
\end{equation}
The right side lies between $A+\E(P-a)^2+\delta^2$ and $A+2\E(P-a)^2+2\delta^2$. When $a$ is an unbiased nominal decoder, the first bound is an equality.
\end{lemma}
\begin{proof}
Since $\E X=1/2$, the oracle Bernoulli variance $\E[(X+\theta)(1-X-\theta)]$ equals $b_\delta$ for either sign. Projection first on the latent score and then on $\mathcal H$ gives $b_\delta+A+\E(P+\theta-a)^2$. Expand the square and maximize over the two signs to obtain \eqref{eq:calibration}. The upper follows from $\abs{\E(P-a)}\le\norm{P-a}_{L^2}$ and $2\delta e\le\delta^2+e^2$. A zero mean error removes the cross term. The conclusion does not treat biased deterministic numerical decoding as automatically unbiased.
\end{proof}

These propositions isolate different lower mechanisms. Acquired small-ball mass constrains label or output resolution. The conditional variance $A$ constrains what the finite acquisition has learned. An inaccessible calibration sign constrains one specified uncertainty class. No statement here asserts a universal positive lower floor for an arbitrary calibration certificate, numerical error bound, or simulation deficiency: an upper allowance can overestimate the error of a particular experiment.

\section{Intermittently mixing finite-state filters}\label{sec:filter}
Fix an integer $d\ge1$, $0<p<1$, $0<a<1$, and $0<s\le1/(d+1)$. Prepare $Z_0$ uniformly on $\{0,1,\ldots,d\}$. At call $t$, independently choose a gate $G_t\sim\operatorname{Bernoulli}(p)$. If $G_t=1$, apply any declared stochastic transition matrix $T_t$ whose entries are at least $s$. If $G_t=0$, leave the hidden state unchanged. The matrices may depend on time and need not have a common invariant law. Conditional on $Z_t=i$, report $V_t\in[-1,1]^d$ with densities
\begin{equation}\label{eq:filter-density}
 g_0(v)=2^{-d},\qquad g_i(v)=2^{-d}(1+av_i),\quad 1\le i\le d.
\end{equation}
The visible report is $(G_t,V_t)$. These are positive normalized densities. The preparation, each call, and the final probe each cost one raw acquisition, so $N_{\rm raw}=T+2$.

Let $\pi_{t,i}=\Pp(Z_t=i)$ be the known marginal obtained from the declared raw protocol. A terminal probe is available at each cutoff; the scored exploration invokes it after $T$ calls. It chooses an index $J$ uniformly from $\{1,\ldots,d\}$ and returns that index and a Bernoulli report of mean
\begin{equation}\label{eq:filter-task}
 X_J+\theta,\qquad X_i=\tfrac12+\tfrac14(1_{\{Z_T=i\}}-\pi_{T,i}),
 \qquad \theta\in\{-\delta,\delta\},\quad 0\le\delta\le1/8.
\end{equation}
No acquisition depends on the unknown terminal sign. A retained label determines a forecast for each possible index; the squared loss averaged over the executed index is the norm $\norm{z}_d^2=d^{-1}\sum_{i=1}^d z_i^2$. This is one actual terminal experiment, not a postulated joint observation of incompatible probes.

The gate-zero steps are not holds. Their observations are informative, but have no hidden mixing. A run of $k$ such steps has probability $(1-p)^k>0$ and is noncontracting in the stable metric below. Thus there is no uniform strict contraction over every fixed-length realized block. The model class has arbitrary finite hidden-state dimension and time-varying positive mixing matrices, rather than one selected scalar refresh model.

\begin{lemma}[Predictive quotient and raw inequalities]\label{lem:filter-raw}
Let $P_{t,i}=\Pp(Z_t=i\mid\cH_t)$ and $L_{t,i}=\log(P_{t,i}/P_{t,0})$ for $1\le i\le d$. At every finite history the posterior is strictly positive and realizes the predictive quotient, with the time/continuation interface included. On $\R^d$ use
\begin{equation}\label{eq:osc-metric}
 D(l,z)=\operatorname{osc}(l_1-z_1,\ldots,l_d-z_d,0),
 \qquad \operatorname{osc}(u)=\max u-\min u.
\end{equation}
Write $P(l)$ for the softmax posterior with reference coordinate zero, and $H_t(l)$ for the log odds of $P(l)T_t$. The update is
\begin{equation}\label{eq:filter-update}
 F_g(l,v)=H_{t,g}(l)+(\log(1+av_i))_{i=1}^d,
 \qquad H_{t,0}(l)=l,\quad H_{t,1}(l)=H_t(l).
\end{equation}
Set $\rho=1-s$, $A_s=\log(1/s)$, and $B_a=\log((1+a)/(1-a))$. Then \eqref{eq:pair} holds with
\[
 \kappa_f=1-p+p\rho^2<1.
\]
Choose $0<\eta< -\log(1-p)/(2B_a)$ and $w(l)=\exp(\eta D(l,0))$. The true-law envelope inequality \eqref{eq:envelope} holds with
\begin{equation}\label{eq:filter-drift}
 a_f=(1-p)e^{2\eta B_a}<1,
 \qquad b_f=p e^{2\eta(A_s+B_a)}.
\end{equation}
These constants are uniform over all the declared matrices and all true and candidate states.
\end{lemma}
\begin{proof}
The raw transition predicts $r=P(l)$ on gate zero and $r=P(l)T_t$ on gate one. Bayes' formula with \eqref{eq:filter-density} gives
\begin{equation}\label{eq:filter-bayes}
 P'_{i}=\frac{r_i(1+av_i)}{D_v},\quad
 P'_0=\frac{r_0}{D_v},\qquad
 D_v=1+a\sum_{i=1}^d r_i v_i.
\end{equation}
It gives \eqref{eq:filter-update} pointwise on every report in the cube. Positivity is preserved. The posterior determines every legal continuation by kernel integration; the indexed current terminal probe separates every one of its coordinates. Hence it realizes the quotient directly, even when a particular mixing matrix has deficient rank.

We prove the needed contraction instead of assuming it for a Bayes recursion. Give the reference log coordinate the value zero, and differentiate the unnormalized predicted log coordinates in a direction $u\in\R^{d+1}$. Their derivatives are expectations of $u$ under the probability vectors
\[
 w_{jk}=\frac{P_j(l)T_{t,jk}}{\sum_hP_h(l)T_{t,hk}}.
\]
Since $T_{t,jk}\ge s$ and the denominator is at most one, $w_{jk}\ge sP_j(l)$ for every output coordinate $k$. The distributions for two output coordinates therefore share the subprobability vector $sP(l)$. Their expectations of $u$ differ by at most $(1-s)\operatorname{osc}(u)$. Integrate the derivative along the line segment between two log-odds vectors to obtain
\[
 D(H_t(l),H_t(z))\le\rho D(l,z).
\]
Every predicted coordinate lies in $[s,1]$, so $D(H_t(z),0)\le A_s$. The additive observation vector cancels from the difference of two updates on the same report. The gate is independent of the entire observed past and algorithm randomization; integration over it gives $\kappa_f$, regardless of the observation law's dependence on the true state.

The oscillation of the observation vector, including its reference zero, is at most $B_a$. Thus $D(F_0(z,v),0)\le D(z,0)+B_a$ and $D(F_1(z,v),0)\le A_s+B_a$. Exponentiate and integrate under any true report law to obtain \eqref{eq:filter-drift}. The drift controls a wrong candidate driven by physical data; it is not a drift under the candidate's imagined observation law.
\end{proof}

\begin{lemma}[A multidimensional inward compander]\label{lem:compander}
For $K\ge1$, put
\begin{equation}\label{eq:compander}
 c_j=\eta^{-1}\log\frac K{K-j},\qquad j=0,\ldots,K-1.
\end{equation}
For $x\ge0$ choose $j(x)=\min\{K-1,\lfloor K(1-e^{-\eta x})\rfloor\}$, set $Q_K(x)=c_{j(x)}$, and extend oddly. Apply this scalar map to each coordinate of $l\in\R^d$. The product cover has $(2K-1)^d$ representatives and obeys
\begin{equation}\label{eq:compander-bound}
 D(Q_Kl,0)\le D(l,0),\qquad
 D(l,Q_Kl)\le\frac{2w(l)}{\eta K}.
\end{equation}
\end{lemma}
\begin{proof}
For a nonnegative scalar $x$, set $u=e^{-\eta x}$ and $v=(K-j(x))/K$. Then $u\le v$, $v-u\le1/K$, also in the unbounded last cell, and
\[
 0\le x-Q_K(x)=\eta^{-1}\log(v/u)
 \le\eta^{-1}(v-u)/u\le e^{\eta x}/(\eta K).
\]
The negative case is symmetric. Each coordinate preserves its sign and decreases its magnitude, so the maximum including reference zero cannot increase and the minimum cannot decrease. This proves inwardness for oscillation. Each scalar error is at most $e^{\eta|l_i|}/(\eta K)\le w(l)/(\eta K)$; their oscillation, including zero, is at most twice their largest magnitude. The product partition is Borel and covers the entire unbounded log-odds domain. There is no truncation of the actual law.
\end{proof}

\begin{theorem}[A class law under intermittent mixing]\label{thm:filter}
In this finite-state experiment, fix $T\ge1$, $M\ge1$, and $b\ge1$. Each physical forecast at the terminal probe must lie in the fixed dyadic grid $\mathcal D_b=\{j2^{-b}:0\le j\le2^b\}$. Let $\mathcal R_\cp$ and $\mathcal R_\on$ be the two sign-minimax risks. Define
\begin{align}
 A_T&=\frac1{16d}\sum_{i=1}^d\E[P_{T,i}(1-P_{T,i})],\label{eq:filter-acquisition}\\
 b_\delta&=\frac14-\frac1{16d}\sum_{i=1}^d
                 \pi_{T,i}(1-\pi_{T,i})-\delta^2.\nonumber
\end{align}
There are $0<c<C<\infty$ depending only on $d,p,s,a$ such that, uniformly over the declared transition matrices, horizon, and budgets,
\begin{equation}\label{eq:filter-law}
 c\{A_T+M^{-2/d}+2^{-2b}+\delta^2\}
 \le\mathcal R_\cp-b_\delta
 \le\mathcal R_\on-b_\delta
 \le C\{A_T+M^{-2/d}+2^{-2b}+\delta^2\}.
\end{equation}
The excess over the sign-aware full-history baseline $b_\delta+A_T$ has matching order $M^{-2/d}+2^{-2b}+\delta^2$. The dimension $d$ is verified by actual acquired posterior mass, not inserted as an ambient parameter count. The online upper follows from Theorem~\ref{thm:inward} and the static product cover, using no retained observation window or real posterior register.
\end{theorem}
\begin{proof}
Let $\zeta_{T,i}=1/2+(P_{T,i}-\pi_{T,i})/4$. Each latent nominal score $X_i$ has expectation $1/2$, and the difference between its physical-oracle and full-history conditional means has squared error $\E[P_{T,i}(1-P_{T,i})]/16$. Average conditional projection over the independent executed probe index. This gives the stated $A_T$ and the same oracle variance $b_\delta$ under both calibration signs. Expanding the two sign risks gives the vector version of Lemma~\ref{lem:calibration}, with cross term
\[
 2\delta\left|\frac1d\sum_{i=1}^d\E(\zeta_{T,i}-a_i)\right|
 \le 2\delta\norm{\zeta_T-a}_{L^2(\norm{\cdot}_d)}.
\]
Thus lower and upper bounds differ by at most the factor-two estimate in that lemma.

Choose the largest $K$ with $(2K-1)^d\le M$. Then $K\ge\tfrac14M^{1/d}$, including small budgets. The known uniform preparation has log odds zero and is exactly a representative. Lemmas~\ref{lem:filter-raw} and \ref{lem:compander} verify Theorem~\ref{thm:inward} with $H_0=1$, $e_0=0$, and $r_M=2/(\eta K)$. The softmax map is Lipschitz from oscillation to each coordinate: its directional derivative is $P_i(u_i-\sum_jP_ju_j)$, of magnitude at most $\operatorname{osc}(u)$. Hence the nominal score is $1/4$-Lipschitz into $\norm{\cdot}_d$. The online score distortion is at most $C M^{-2/d}$. Rounding each possible probe forecast to $\mathcal D_b$ adds a root error at most $2^{-b}$. Projection and the sign expansion give the upper in \eqref{eq:filter-law}.

For a lower use only the actual event $G_T=1$, of mass $p$. Conditional on any previous history and this event, every predicted coordinate satisfies $r_i\ge s$. The map in \eqref{eq:filter-bayes} from the report cube to $(P'_1,\ldots,P'_d)$ is injective, with inverse
\[
 v_i=\frac1a\left(\frac{r_0P'_i}{r_iP'_0}-1\right).
\]
Its Jacobian determinant, by the rank-one determinant formula, is
\begin{equation}\label{eq:filter-jacobian}
 \left|\det\frac{\partial(P'_1,\ldots,P'_d)}{\partial(v_1,\ldots,v_d)}\right|
 =\frac{a^d\prod_{i=0}^d r_i}{D_v^{d+1}}
 \ge \frac{a^d s^{d+1}}{(1+a)^{d+1}}>0.
\end{equation}
The report density is $2^{-d}D_v\le2^{-d}(1+a)$. Change of variables bounds the conditional posterior density by a constant depending only on $d,s,a$; scaling each score coordinate by $1/4$ preserves such a bound with another constant $C_*$. Mixing the conditional densities over previous histories preserves it. The actual submeasure $\underline\nu_T(E)=\Pp(\zeta_T\in E,G_T=1)$ therefore has mass $p$ and, in the task norm, obeys
\[
 \sup_u\underline\nu_T(B(u,r))
 \le pC_*v_d d^{d/2}r^d,
\]
where $v_d$ is the Euclidean unit-ball volume. Proposition~\ref{prop:submass} gives $q_m(\law(\zeta_T))\ge c_fm^{-2/d}$ for all $m\ge1$. No bounded density is assumed for the gate-zero part of the acquired law.

A label determines a vector of $d$ possible probe forecasts. Under the output restriction this vector belongs to $\mathcal D_b^d$, of cardinality $(2^b+1)^d$. Proposition~\ref{prop:output} and the average of the two sign risks yield the lower
\[
 A_T+\delta^2+q_{\min(M,(2^b+1)^d)}(\law(\zeta_T)).
\]
The quantization term is at least a constant times $M^{-2/d}+2^{-2b}$. This proves \eqref{eq:filter-law} and the same comparison after subtracting the common acquisition term. The quantifiers include every finite dimension $d$, with dimension-dependent constants explicitly allowed.
\end{proof}

\subsection{Finite input and arithmetic}
The theorem uses a continuous report as an exact Borel input, not a free finite communication packet. A certified box of coordinate widths at most $h$ containing the exact candidate can be shrunk coordinatewise toward zero: choose zero if the interval contains it, the left endpoint if positive, and the right endpoint if negative. The result has no larger oscillation including zero and is within $2h$ in the metric $D$. Corollary~\ref{cor:perturb} therefore applies with $u=2h$.

For the observation log likelihood, the sensitivities to a report coordinate and to $a$ are bounded by $a/(1-a)$ and $1/(1-a)$. For two admissible mixing matrices differing entrywise by at most $e$, every predicted coordinate differs by at most $e$; if both matrices have entries at least $s$, their predicted log coordinates differ by at most $e/s$, giving an oscillation error at most $2e/s$. The identity branch is implemented exactly. Certified log-sum-exp and likelihood evaluations thus produce a common enclosure from separate matrix-calibration, input, and arithmetic allowances. The nominal marginal $\pi_T$ used in the readout is computed from the declared protocol; any error in it is charged in the readout allowance $v$.

Round the positive scalar table nodes downward by at most $\tau$ and select the largest rounded node not exceeding a candidate coordinate's magnitude. They remain inward. The scalar covering error increases by at most $\tau$, so the product is an inward cover of radius $2/(\eta K)+2\tau$. All comparisons are rational and terminate, including ties. For computable parameters, interval logarithm routines provide any fixed table tolerance. Table descriptions or on-demand generation, temporary arithmetic, input packets of $d$ coordinates, and the controller for the matrix schedule are separate program and workspace costs. The persistent prediction state is only the product-table index.

The constants deteriorate as mixing vanishes or as the acquired observation loses rank. Informative gate-zero calls cannot be relabeled as exact holds to remove their genuine dependence on $p$. Nor is $A_T$ asserted to have an $N^{-1}$ rate: the hidden dynamics continually create uncertainty. Kernel learning, optimized interventions, and stopping-dependent acquisition need additional theorems rather than reinterpretation of this fixed-protocol result.

\section{Finite acquisition through nonreset singular dynamics}\label{sec:singular}
We next verify the same raw-kernel theorem on a compact singular domain with expanding updates. The construction extends the non-erasing symbolic realization in the supplement, but does not require its invariant-reference preparation condition. The initial mode law and the transition probabilities below may be time dependent or have zero coordinates.

Fix ratios $r_n\in[1/5,1/3]$, let $\ell_0=1$, $\ell_n=\prod_{j=1}^n r_j$, and encode a binary sequence $\omega$ by
\begin{equation}\label{eq:cantor}
 x(\omega)=\sum_{n\ge1}(1-r_n)\ell_{n-1}\omega_n.
\end{equation}
The image $K$ is a compact Cantor set. For a finite word $u$ of length $h$, write
\begin{equation}\label{eq:cylinder}
 c(u)=\sum_{n=1}^h(1-r_n)\ell_{n-1}u_n+\tfrac12\ell_h,
 \qquad c(\varnothing)=\tfrac12.
\end{equation}
This is the mean of the fair-tail posterior in that cylinder. Let $K^*$ consist of $K$ and all the finite-word centers. This is a compact subset of $[0,1]$: a sequence of centers with unbounded lengths approaches $K$, and the centers of bounded length form a finite set. A terminating center is a predictive state, not a physical point of $K$.

Prepare a mode $I_0\in\{0,1\}$ with any declared law, independently of a fair infinite tail. The preparation reports only the mode. Make $B$ paid calls, each reporting the next initial tail bit. Thus the full-history state after these calls is $(I_0,c(U_B))$, where $U_B$ is the acquired word. There is no infinite initialization report.

There are then $T$ recurrent opportunities. At time $t$, an independent active indicator has arbitrary deterministic probability $\vartheta_t\in[0,1]$. On a hold the physical mode and tail, and the retained state, are copied exactly. On an active step choose the following operation, with $\alpha_t\in[0,1]$ arbitrary and deterministic, and $\gamma=1/4096$:
\begin{center}
\begin{tabular}{ccll}
\toprule
Source & Destination & Probability within the source & Tail operation \\
\midrule
$0$ & $0$ & $(1-\alpha_t)(1-\gamma)$ & prepend two fair bits \\
$0$ & $0$ & $(1-\alpha_t)\gamma$ & delete first bit \\
$0$ & $1$ & $\alpha_t$ & delete first bit \\
$1$ & $0$ & $1/2$ & prepend six fair bits \\
$1$ & $1$ & $1/2$ & copy \\
\bottomrule
\end{tabular}
\end{center}
Report the destination, operation, and newly prepended bits. Do not report the deleted bit or the surviving old tail. Both changes of mode retain a nonconstant old tail. These operations define nonnegative normalized Borel kernels, and every bit in a new report has its specified fair law. Mode and operation probabilities are independent of tail-bit values.

Let $w_T=\Pp(I_T=1)$ under this \emph{actual} finite protocol. It is a known function of the declared mode law and schedule, not of the unknown terminal sign. The terminal probe has mean
\begin{equation}\label{eq:symbol-score}
 X+\theta,\qquad
 X=\tfrac12+\tfrac14(I_T-w_T)+\tfrac18(x(\omega_T)-\tfrac12),
 \quad \theta\in\{-\delta,\delta\},\quad0\le\delta\le1/16.
\end{equation}
These means lie in $[0,1]$. The fair-tail marginal is preserved by every edge, so $\E X=1/2$, even for nonstationary mode laws. The score intervals of the two modes are disjoint. Within a mode the nominal terminal score is strictly affine in the cylinder mean. Hence the mode and cylinder posterior realize the predictive quotient of the declared finite protocol. They determine all later reports, and the terminal probe separates them.

\begin{lemma}[Completed-domain gains and inward covers]\label{lem:symbol-gains}
On $K^*$, deletion (fixing $c(\varnothing)$) has Lipschitz constant at most $30$, prepending a fixed word of length $s$ has constant at most $6\,3^{-s}$, and copying has constant one. For each $m\ge0$, use all centers of words of lengths at most $m$ in both modes, a total of
\begin{equation}\label{eq:symbol-count}
 M_m=2^{m+2}-2
\end{equation}
labels. Leave a shorter center fixed, and send a longer center or an infinite sequence to the center of its length-$m$ prefix. On
\[
 S=\{0,1\}\times K^*,\qquad
 d((i,x),(i,y))=\sqrt{v_i}|x-y|,
 \quad v_0=1800,
 \quad v_1=1,
\]
with cross-mode distance $2\sqrt{1800}$, this is a Borel inward cover for the constant envelope $w=1$ and radius $\sqrt{1800}\,\ell_m$.
\end{lemma}
\begin{proof}
A finite center can be viewed as a terminated word, with the remainder replaced by its fair mean. If two completed coordinates first differ or one terminates at level $k$, their distance is between $\ell_{k-1}/6$ and $\ell_{k-1}$. Indeed, opposite binary branches have gap at least $(1-2r_k)\ell_{k-1}\ge\ell_{k-1}/3$; a center at the parent mean and either branch are separated by at least $(1/2-r_k)\ell_{k-1}\ge\ell_{k-1}/6$. This also proves unique decoding of every terminating center and infinite point.

If two nonterminated first digits agree, deletion shifts the first distinguishing level down by one. The ratio between the new upper gap and old lower gap is at most $6/r_{k-1}\le30$. If the first digits differ, or one coordinate is the empty-word mean, the old gap is at least $1/6$ and the new gap at most one; this is smaller than the same bound. Prepending $s$ common digits shifts the distinguishing level up by $s$, giving a ratio at most $6\ell_{k+s-1}/\ell_{k-1}\le6\,3^{-s}$. Copying is immediate.

The stated distances form a metric: each within-mode diameter is at most $\sqrt{1800}$, while every cross-mode distance is $2\sqrt{1800}$. Prefix cylinders and individual terminating centers are Borel. The representative error within a length-$m$ cylinder is at most $\ell_m$, and shorter centers are unchanged. The mode is part of the label. Summing $2^h$ over $0\le h\le m$ and multiplying by two gives \eqref{eq:symbol-count}. With $w=1$, the inward inequality is automatic.
\end{proof}

\begin{lemma}[Raw two-point drift without occupation comparison]\label{lem:symbol-drift}
The active report kernel on $S$ satisfies \eqref{eq:pair} with $\kappa=17/20$, uniformly in all $\alpha_t\in[0,1]$. It satisfies \eqref{eq:envelope} with $w=1$, $a=0$, and $b=1$. These conclusions do not depend on the preparation or any terminal mode mass.
\end{lemma}
\begin{proof}
For a candidate in a different mode from the true state, extend the update by applying the \emph{reported} tail operation to its coordinate and setting its mode to the reported destination. This is a Borel map declared on all candidate inputs. It uses no inaccessible information. For a same-mode pair, Lemma~\ref{lem:symbol-gains} bounds squared gains through the matrix
\begin{equation}\label{eq:symbol-matrix}
 C_t=\begin{pmatrix}
 (1-\alpha_t)85/128 & 900\alpha_t\\
 2/59049 & 1/2
 \end{pmatrix}.
\end{equation}
The first entry uses
\[
 (1-\gamma)\frac49+900\gamma=\frac{85}{128}.
\]
For $v=(1800,1)^{\mathsf T}$, the first component of $C_tv$ is at most $1800(85/128)<(17/20)1800$, and the second is $3600/59049+1/2<17/20$. Consequently the expected squared destination-weighted gap is at most $(17/20)$ times the source-weighted gap. For a cross-mode pair, the old squared distance is $4\cdot1800$ and after every report both modes agree, with new squared distance at most $1800$. Its gain is at most $1/4$. This proves \eqref{eq:pair} for every pair in $S$, not only pairs already reached by a particular code. The envelope equality is immediate. The entire calculation is under each true state's raw report law and retains both off-diagonal transitions.
\end{proof}

\begin{theorem}[Actual-depth resource law with changing mode masses]\label{thm:singular}
In the finite-read experiment above, let $H_t$ denote the length of the prefix known to the full history after $t$ recurrent opportunities. Start with $H_0=B$ and update it by
\[
 h\mapsto h+2,\quad h\mapsto h+6,\quad
 h\mapsto\max(h-1,0),\quad h\mapsto h
\]
under the respective reported operations. Let $\omega_{T,i,h}=\Pp(I_T=i,H_T=h)$ be the actual probabilities obtained by this finite recursion. For the nominal conditional score law,
\begin{equation}\label{eq:atomic-law}
 \nu_{B,T}=\sum_{i,h}\omega_{T,i,h}\,2^{-h}
      \sum_{u\in\{0,1\}^h}
       \delta_{\,1/2+(i-w_T)/4+(c(u)-1/2)/8}.
\end{equation}
Thus $\nu_{B,T}$ is finite atomic, supported on acquired cylinder means, not on an assumed continuous latent geometry.

Put $b_\delta=\E[X(1-X)]-\delta^2$ for the common physical-oracle baseline in \eqref{eq:symbol-score}, and
\[
 A_{B,T}=\E\bigl(X-\E[X\mid\cH_{B,T}]\bigr)^2.
\]
For every $M\ge6$, writing $n=\lfloor\log_2M\rfloor$, the sign-minimax checkpoint and online risks satisfy
\begin{align}
 \mathcal R_\cp-b_\delta
    &=\delta^2+A_{B,T}+q_M(\nu_{B,T}),\label{eq:symbol-exact}\\
 \mathcal R_\cp-b_\delta
    &\asymp\mathcal R_\on-b_\delta
     \asymp\delta^2+\E\ell_{\min(H_T,n)}^2.\label{eq:symbol-law}
\end{align}
The comparison constants are universal over the ratios, $B,T$, the initial mode law, the schedules $(\alpha_t)$ and $(\vartheta_t)$, and the allowed signs. The total raw acquisition count is $B+T+2$. An implementing encoder stores only a mode and a word of bounded length, including its length in the label count \eqref{eq:symbol-count}.
\end{theorem}
\begin{proof}
Initially the acquired word is uniform, and the unobserved remainder is fair and independent of it and the mode. Prepending reveals fresh independent fair bits, deleting removes the first bit without revealing an unobserved one, and copying changes nothing. Induction over the reports shows that the full history always knows exactly the stated prefix and that its unobserved remainder is fair. The depth and mode evolution depend on operations, not bit values; conditioning only on $(I_T,H_T)=(i,h)$ therefore leaves the known word uniform. This proves the actual atomic law \eqref{eq:atomic-law}. Its support is finite since $H_T\le B+6T$.

Within a depth-$h$ cylinder, the physical coordinate variance is between $\ell_h^2/9$ and $\ell_h^2/4$. The first remaining independent fair bit has coefficient $(1-r_{h+1})\ell_h\ge(2/3)\ell_h$, which proves the lower; a variable in an interval of length $\ell_h$ has variance at most $\ell_h^2/4$. The score slope $1/8$ consequently gives
\begin{equation}\label{eq:symbol-acquisition}
 \frac1{576}\E\ell_{H_T}^2
 \le A_{B,T}\le\frac1{256}\E\ell_{H_T}^2.
\end{equation}

For the online construction choose $m=\lfloor\log_2(M+2)\rfloor-2\ge1$ and the static cover in Lemma~\ref{lem:symbol-gains}. The $B$ initial readings are processed once: retain the first $\min(B,m)$ bits and discard any others. This is a legal seed with error at most $\sqrt{1800}\ell_m$ and envelope one. No full-history checkpoint routine is used for initialization. On each later report apply the raw completed-domain update and then the static cover. Lemma~\ref{lem:symbol-drift} and Theorem~\ref{thm:inward} give nominal acquired-score distortion at most $C\ell_m^2\le C'\ell_n^2$, since $0\le n-m\le2$ and every ratio is at least $1/5$. The score is Lipschitz for the stated metric, including its cross-mode distance.

Equivalently, the resulting concrete machine prepends and truncates a retained word, deletes its first bit if present, or copies it. Its retained word is always an actual prefix of the physical tail. Conditional on the origin pattern of the retained bits and on the retained word, the unretained bits are independent fair bits. The origin pattern depends only on the operation history, not on bit values. Averaging over this pattern shows that the nominal center readout is the conditional mean of $X$ given the retained mode and word. It is therefore unbiased. Lemma~\ref{lem:calibration} adds exactly $\delta^2$ for this readout and yields the upper $\delta^2+A_{B,T}+C'\ell_n^2$.

For the lower use the \emph{physical} nominal score only through the common-oracle inequality, Proposition~\ref{prop:oracle}. Every operation preserves a fair physical tail, so conditional on each terminal mode it has the Cantor law from \eqref{eq:cantor}, whatever the initial mode law and schedule. At least one terminal mode has mass at least $1/2$. At level $k$, cylinders have diameter $\ell_k$ and gaps at least $\ell_k$. A score ball of radius $\ell_k/32$ meets at most one cylinder of this chosen mode, because its score diameter is $\ell_k/16$ whereas the within-mode score gaps are at least $\ell_k/8$. The conditional cylinder mass is $2^{-k}$. Set $k=n+2$; then $2^{-k}\le1/(2M)$ and $\ell_k\ge\ell_n/25$. The acquired-submeasure union argument applied to the ideal score law therefore gives
\begin{equation}\label{eq:ideal-cantor-lower}
 q_M(\law(X))\ge c\ell_n^2
\end{equation}
with a universal $c>0$. This is not a small-ball assertion about the atomic law \eqref{eq:atomic-law}.

By averaging the two sign risks, every checkpoint code has excess at least $\delta^2+A_{B,T}+q_M(\nu_{B,T})$. Conversely an optimal nominal checkpoint code exists for the finite atomic law. Replace each of its readouts by the conditional mean of $X$ given that label; this cannot increase nominal risk and keeps at most $M$ labels. Its mean error is zero, so Lemma~\ref{lem:calibration} attains that same lower bound. This proves \eqref{eq:symbol-exact}. Equations \eqref{eq:symbol-acquisition} and \eqref{eq:ideal-cantor-lower}, the common-oracle identity, and the online upper imply comparison with $\delta^2+\E\ell_{H_T}^2+\ell_n^2$. Finally monotonicity of $\ell_h$ gives
\[
 \max\{\E\ell_{H_T}^2,\ell_n^2\}
 \le\E\ell_{\min(H_T,n)}^2
 \le\E\ell_{H_T}^2+\ell_n^2,
\]
which proves \eqref{eq:symbol-law}. All waits and initial reads are charged in $B+T+2$.
\end{proof}

\begin{corollary}[A singular profile without a single exponent]\label{cor:no-exponent}
There are ratios $r_n\in\{1/3,1/5\}$ for which the geometric factor $\ell_{\lfloor\log_2M\rfloor}^2$ has two distinct subsequential values of its negative logarithmic slope in $M$, namely $2\log3/\log2$ and $2\log5/\log2$. The same slopes occur for the risk in \eqref{eq:symbol-law} along the finite-acquisition protocols $T=0$, $\delta=0$, $B=\lfloor\log_2M\rfloor+2$.
\end{corollary}
\begin{proof}
Choose alternating constant-ratio blocks whose lengths exceed the square of the sum of all preceding lengths. At the end of each block, the Ces\`aro average of $-\log r_j$ converges along the corresponding subsequence to $\log3$ or $\log5$. Since $\log M/\lfloor\log_2M\rfloor\to\log2$, the geometric slopes follow. For the stated finite protocols $H_T=B\ge n$, so \eqref{eq:symbol-law} reduces to $\ell_n^2$ up to universal constants. This is a family of finite acquisition experiments with increasing $B$; it is not a claim of a nonzero quantization exponent for a fixed finite atomic law below its acquisition floor.
\end{proof}

Only a finite packet is read at a recurrent call: a mode, an operation, and at most six fresh bits. The clock needed to follow a nonconstant schedule, the schedule's program description, and arithmetic for the finitely many ratios remain separate accounts. For rational ratios $1/3$ or $1/5$, a length-$m$ center is rational with numerator and denominator of $O(m)$ bits. Straightforward integer arithmetic gives a finite implementation with $O(m+b)$ temporary bits for a $b$-bit output approximation, in addition to the schedule interface. The statement is an implementation upper, not an optimal time or workspace lower. The exact loss identity \eqref{eq:symbol-exact} uses the unbiased nominal centers; arbitrary finite rounding must instead use the cross-term bound in Lemma~\ref{lem:calibration}.

\section{Causal morphisms and separated resource accounts}\label{sec:resources}
A resource statement is meaningful only after the causal information cut is fixed. In particular, the predictive label count is not a surrogate for all writable bits or for an uncharged simulator. We give the common-loss composition explicitly, and then state the exact additional condition under which an interface factor really forces more retained states.

\subsection{Report continuity}
\begin{proposition}[A finite-horizon report comparison]\label{prop:report}
In Theorem~\ref{thm:inward}, suppose the active report kernels obey
\[
 \TV(J_t(\cdot\mid x),J_t(\cdot\mid z))\le L_Jd(x,z).
\]
Suppose also that the terminal task-report kernel obeys the analogous bound with constant $L_Y$. The nominal exact experiment and the recursively quantized experiment that generates reports from the current representative, using the same legal controller and independent holds, have complete report-and-task laws at total variation distance at most
\begin{equation}\label{eq:path-tv}
 \min\{1,(TL_J+L_Y)E_M\}.
\end{equation}
The assertion is for the prescribed finite horizon, with controller randomness coupled identically. Constants must hold uniformly over that controller's declared actions. It is not a comparison of separately optimized unrestricted controllers.
\end{proposition}
\begin{proof}
Construct successive maximal couplings of the two report kernels until their first disagreement. Standard Borel report spaces admit measurable such couplings, for example by taking the common part of the two kernels relative to their sum and splitting the remainders. Before the first disagreement, both controllers see the same reports and randomization and therefore choose the same action. The representative process equals the recursion driven by the true reports. Define that shadow recursion also after a disagreement, so that Lemma~\ref{lem:recursion} always applies to it under the physical law. The probability of a first disagreement at time $t$ is at most the expectation of $L_Jd(S_{t-1},\widehat S_{t-1})$ on the event of no previous disagreement, hence at most $L_JE_M$. Exact holds create no such error. Couple the terminal task similarly at cost at most $L_YE_M$. The union bound gives \eqref{eq:path-tv}. Total variation is here $\sup_A|P(A)-Q(A)|$, so no extra factor is used for bounded losses.
\end{proof}

For the filter, the report density given a predicted probability vector $r$ is $2^{-d}(1+a\sum_i r_i v_i)$. Its total variation difference is at most $(a/4)\sum_{i=1}^d|r_i-r'_i|$. Differentiating softmax along a log-odds segment bounds the full $\ell^1$ probability difference by its oscillation distance. The mixing prediction is nonexpansive in that metric, so $L_J=a/4$ is valid. The indexed nominal terminal probe has a valid constant $L_Y=1/4$. For the singular experiment the active report law is independent of the coordinate within a mode; across modes its total variation distance is at most one and the declared state distance is $2\sqrt{1800}$. Both realizations therefore verify report continuity directly from their raw kernels. The bound \eqref{eq:path-tv} may grow with the horizon even when prediction risk is uniformly bounded; uniform task accuracy is not infinite-horizon path-law equivalence or large-deviation equivalence.

\subsection{A typed composition theorem}
Let $\mathfrak E=(P_\lambda)_\lambda$ and $\mathfrak F=(Q_\lambda)_\lambda$ be parameter-indexed prepared causal experiments. A morphism from $\mathfrak E$ to $\mathfrak F$ specifies a parameter-independent causal simulator, a transformation of reports, a lift of each declared target strategy to legal source actions, an internal simulator state, a physical-to-virtual clock, and compatible terminal task and action readouts. Each strategy lift includes its stopping and action-availability conventions. The simulator is not allowed to read the unknown $\lambda$.

Its complete interactive defect is a bound $\varepsilon$ on the total variation distance between the virtual target transcript together with the common scored outcome, and the true target law, uniformly in $\lambda$ and in the declared strategy class. This is stronger than a one-step marginal or fixed-controller terminal comparison. An implementation specifies bounds for raw calls $N$, physical time $\tau$, prediction labels $M$, simulator states $R$, controller states $D$, phase/clock states $Q$, temporary workspace $W$, and read-only program description $L_{\rm prog}$. Values may depend on the declared horizon and accuracy. Infinite or unspecified overhead does not yield a finite total-state theorem.

\begin{theorem}[Common-risk and resource composition]\label{thm:morphism}
Suppose a morphism as above simulates a target policy with risk $r_\lambda$ for a common loss in $[0,L_*]$. Suppose its target predictor has $M$ labels, its simulator has $R$ states, its lifted controller has $D$ states, and its phase/clock has $Q$ states, with all bounds on the complete legal execution. Then there is a serial source implementation with
\begin{equation}\label{eq:resource-product}
 K\le MRDQ
\end{equation}
retained states and risk at most $r_\lambda+L_*\varepsilon$ for every $\lambda$. Raw calls and time add over the executed subroutines; simultaneously needed workspace and stored program descriptions have their corresponding additive upper accounts.

If a second such morphism is composable and the first lifted strategies remain in its certified strategy class, the complete defects add, capped at one. The retained-state overheads multiply for the serial product implementation. In particular, applying Theorem~\ref{thm:inward} and Corollary~\ref{cor:perturb} in the target gives the common-risk upper
\begin{equation}\label{eq:full-composition}
 B_T+(L E_{M,u}+v)^2+L_*\varepsilon,
 \qquad K\le MRDQ,
\end{equation}
whenever the target and source scored tasks are the common compatible task just described. The $B_T$ in this expression is the stated target benchmark, not a separately recomputed source Bayes risk.
\end{theorem}
\begin{proof}
At every step retain the tuple of predictor, simulator, controller, and clock states, run the prescribed source subroutine, transform its report, update the target label, and apply the compatible readout. The tuple set has cardinality at most $MRDQ$. Parameter independence, causal strategy lifting, and action compatibility give a legal source procedure. No virtual action is performed before the source reports needed to choose it exist. The definition of complete defect and the bounded-loss total variation inequality give the risk bound. Counting all physical subroutine calls includes failed attempts, calibration calls, and stopping markers. Sequential clocks and costs add, and storing all concurrently needed registers and program descriptions gives the asserted upper accounts.

For composition, first insert the intermediate simulated law, apply the triangle inequality for total variation, and use its contraction under the parameter-independent second causal transformation. The strategy-class closure hypothesis allows the second defect bound to be applied to the actually lifted adaptive policy. Combining the two tuples multiplies their cardinality bounds. For \eqref{eq:full-composition}, the metric approximation terms have already been combined in the target task norm before squaring, and the complete common-loss defect is added afterward. A comparison of different full-history Bayes baselines is neither required nor inferred.
\end{proof}

\subsection{When an interface state is necessary}
The product in \eqref{eq:resource-product} is generally an upper bound only. The following continuation-separation principle, also used in the supplement, states a condition yielding an actual converse.

\begin{proposition}[An executable continuation cut]\label{prop:interface}
Before a memory cut, an experiment reports a tag $J\in\{1,\ldots,D\}$ with probabilities $p_j>0$ and a task score $U$, whose conditional law given $J=j$ is $\nu_j$. After the cut no report reveals $J$ or $U$. A legal continuation can require exact recovery of $J$, and the procedure must pass this challenge with probability one. It also forecasts $U$ with squared loss. A checkpoint encoder retaining at most $K$ states has optimal distortion
\begin{equation}\label{eq:interface-allocation}
 \inf_{\substack{m_j\ge1\ \sum_jm_j\le K}}
       \sum_{j=1}^D p_j q_{m_j}(\nu_j),
\end{equation}
with infeasibility if $K<D$. Randomization independent of $(J,U)$ cannot improve the value. When $J$ is uniform and each $\nu_j$ is uniform on $[0,1]^d$, this has order $(D/K)^{2/d}$ for $K\ge D$, with constants depending only on $d$.
\end{proposition}
\begin{proof}
Fix an independent shared seed outside a null set on which the exact challenge could fail. A retained label cannot have positive conditional probability under two different tags: its tag-readout distribution would then have to equal two distinct point masses. Hence the labels split into disjoint supports, with sizes $m_j\ge1$ and sum at most $K$. Conditional projection and nearest-center quantization under each tag give the lower in \eqref{eq:interface-allocation}. Conversely assign disjoint label lists to the tags, quantize $U$ within each list, and read the tag and center from the list identity. Approximate quantizers approach the displayed infimum. Seed conditioning proves that randomization does not reduce it.

For the uniform cube, a ball of radius $r$ has mass at most $v_dr^d$, giving $q_m\ge c_dm^{-2/d}$ by Proposition~\ref{prop:submass}. A regular grid with $\lfloor m^{1/d}\rfloor^d$ cells gives the upper $C_dm^{-2/d}$. Jensen's inequality supplies $D^{-1}\sum_j m_j^{-2/d}\ge(K/D)^{-2/d}$; assigning each tag at least $\lfloor K/D\rfloor$ labels gives the reverse inequality up to a dimension-dependent constant. Exact tag recovery is an executable requirement, not an assumption that an arbitrary simulator's private state is indispensable.
\end{proof}

If an actual continuation challenges an independent tuple of simulator, controller, and clock tags, Proposition~\ref{prop:interface} can make their joint tag cardinality necessary. In the absence of such a challenge or another lower witness, no necessity of $RDQ$ is concluded. Likewise, the output-alphabet lower in Proposition~\ref{prop:output} is not a lower on all internal floating-point arithmetic. Workspace, program length, computation time, and adaptive acquisition still require task- and machine-specific converse arguments.

\section{Position of the results and remaining mathematical boundaries}\label{sec:comparison}
\subsection{Standard tools and the new construction}
Conditional projection, nearest-center quantization, small-ball converses, and bounded-loss comparison are standard tools, not new foundations in themselves. The quantization background is developed in Graf and Luschgy~\cite{graf}; comparison of experiments in Blackwell~\cite{blackwell}, Le Cam~\cite{lecam}, and Torgersen~\cite{torgersen}. Predictive-state representations and future-equivalence constructions have substantial antecedents in Littman, Sutton, and Singh~\cite{psr}, Shalizi and Crutchfield~\cite{shalizi}, and kernel sufficiency treatments such as Fritz~\cite{fritz}. Proposition~\ref{prop:quotient} records the additional pointwise instrument and measurable-image hypotheses actually used here; it is not a claim to originate predictive sufficiency.

The mathematical contribution of Theorem~\ref{thm:inward} is the constructive implication from a static \emph{inward} cover and true-law raw moment bounds to a legal finite recursive state, with a converse stated in the geometry of the actual experiment. Inwardness controls the approximate trajectory's envelope, while the two-point inequality controls its separation from the true trajectory. Neither an existing recursive quantizer nor intermediate-to-terminal occupation domination is a premise. The proofs use standard $L^2$ and Lyapunov arguments; the claim concerns their specified causal-geometric implication and its verified realizations, not the discovery of these analytic inequalities.

Average contraction of random Lipschitz maps is classical; see Diaconis and Freedman~\cite{diaconis}. Our square-moment condition is stronger than a mere negative logarithmic average and is not a replacement for that theory. It is chosen because it controls accumulated squared task error under report laws that depend on the true predictive state. The additional envelope is needed when a finite code on an unbounded stable coordinate has resolution depending on location. The product compander supplies an explicit construction in that setting. It does not claim that every average-stable random system admits the same rate.

\subsection{Filtering approximation and causal coding}
Quantization methods for nonlinear filtering are developed by Pag\`es and Pham~\cite{pages}, Sellami~\cite{sellami}, and Feuer and Goodwin~\cite{feuer}. The present paper does not claim novelty for recursively quantizing a filter or for propagating approximation error. Its finite-state realization adds a matching converse for the \emph{actually acquired} posterior law, uniform over a declared time-varying class with intermittent mixing, and combines that converse with output-alphabet and calibration restrictions in one scored experiment. The lower uses only the genuine gate-one submeasure and a raw Jacobian calculation, not a global posterior-density hypothesis.

McDonald and Y\"uksel~\cite{mcdonald} establish true-law expected filter contraction through transition and observation Dobrushin coefficients. This is an important antecedent for using the physical report law rather than the incorrect filter's imagined law. Our finite-state calculation uses an elementary common-mass bound in log coordinates and additionally constructs an inward finite-label recursion.

Kara and Y\"uksel~\cite{kara}, in particular their Theorems~12, 16, and 17, relate finite-window control approximation to filter stability under specified discount and regularity conditions. Their control objective is broader than our fixed-exploration prediction objective. Our label budget is the cardinality of a recursive retained state, rather than a prescribed observation-window length; our continuous-report example adds an actual geometric lower bound. Neither statement implies the other without an additional conversion between legal policy classes and resources. Saldi, Y\"uksel, and Linder~\cite{saldi} likewise study near-optimal control through finite belief-model approximations. We do not reinterpret a same-controller comparison as an optimal-control theorem.

Nonanticipative rate-distortion theory, including Charalambous, Stavrou, and Ahmed~\cite{nonanticipative}, imposes causal realizability on reconstruction kernels. Its information objective is not identical to a hard bound on every retained-state alphabet, even though information inequalities can supply converse tools. Proposition~\ref{prop:interface} identifies a precise continuation requirement forcing disjoint label supports. It does not turn every simulator implementation into an information-theoretically necessary product overhead.

\subsection{Singular geometry and complementary results}
The geometry of Moran quantization has an extensive theory; Zhu~\cite{zhu-moran} studies asymptotic uniformity for such measures. The novelty asserted in Theorem~\ref{thm:singular} is not a new Cantor-dimension formula. Its object is a finite, physically acquired cylinder law transported recursively through expanding and non-erasing transitions. The common-oracle identity keeps its atomic acquisition variance, while the static completed-domain cover and raw pair drift construct the online machine without an invariant preparation comparison.

The companion mathematical supplement~\cite{supplement} preserves the earlier allocation-weighted causal resolvent, terminal-bridge, finite-chart, robust-mixture, single-probe, and soft-state information results. These remain complementary: for example, nonuniform component allocations or separately certified continuation-information budgets need not satisfy a uniform inward cover at a single common scale. Conversely the present transient raw-drift argument does not require their occupation-comparison certificate. We claim neither implication without checking the respective hypotheses.

\subsection{Scope}
The general theorem and the two class verifications are complete under their displayed assumptions. Several larger objectives remain distinct proof obligations. A necessary-and-sufficient characterization of when actual checkpoint geometry is recursively realizable is not proved. Nor is a joint optimal region for persistent state, temporary workspace, program description, computation time, unknown-kernel learning, and causal deficiency. The continuation cut and output-alphabet theorem give genuine converses for their specified restrictions, not that entire region. Adaptive exploration and error-dependent stopping can change both the acquired law and the conditional inequalities; they require new analysis rather than a formal minimization of the present fixed-protocol bound.

These are boundaries of the conclusions, not substitutions for the experimental problem. The results establish a general constructive route and two different raw-kernel verifications, while keeping the broader acquired-geometry and resource-transfer question mathematically explicit.

\begin{thebibliography}{99}
\bibitem{blackwell} D. Blackwell, \emph{Equivalent comparisons of experiments}, Ann. Math. Statist. \textbf{24} (1953), 265--272. \href{https://doi.org/10.1214/aoms/1177729032}{doi:10.1214/aoms/1177729032}.
\bibitem{diaconis} P. Diaconis and D. Freedman, \emph{Iterated random functions}, SIAM Rev. \textbf{41} (1999), 45--76. \href{https://doi.org/10.1137/S0036144598338446}{doi:10.1137/S0036144598338446}.
\bibitem{feuer} A. Feuer and G. C. Goodwin, \emph{On-line quantization in nonlinear filtering}, J. Stat. Comput. Simul. \textbf{83} (2013), 1210--1222. \href{https://doi.org/10.1080/00949655.2012.656641}{doi:10.1080/00949655.2012.656641}.
\bibitem{graf} S. Graf and H. Luschgy, \emph{Foundations of Quantization for Probability Distributions}, Lecture Notes in Mathematics 1730, Springer, 2000. \href{https://doi.org/10.1007/BFb0103945}{doi:10.1007/BFb0103945}.
\bibitem{psr} M. L. Littman, R. S. Sutton, and S. Singh, \emph{Predictive representations of state}, Advances in Neural Information Processing Systems \textbf{14} (2001), 1555--1561.
\bibitem{pages} G. Pag\`es and H. Pham, \emph{Optimal quantization methods for nonlinear filtering with discrete-time observations}, Bernoulli \textbf{11} (2005), 893--932. \href{https://doi.org/10.3150/bj/1130077599}{doi:10.3150/bj/1130077599}.
\bibitem{sellami} A. Sellami, \emph{Quantization based filtering method using first order approximation}, SIAM J. Numer. Anal. \textbf{47} (2010), 4711--4734. \href{https://doi.org/10.1137/060652580}{doi:10.1137/060652580}.
\bibitem{shalizi} C. R. Shalizi and J. P. Crutchfield, \emph{Computational mechanics: pattern and prediction, structure and simplicity}, J. Stat. Phys. \textbf{104} (2001), 817--879. \href{https://doi.org/10.1023/A:1010388907793}{doi:10.1023/A:1010388907793}.
\bibitem{saldi} N. Saldi, S. Y\"uksel, and T. Linder, \emph{Finite model approximations for partially observed Markov decision processes with discounted cost}, arXiv version consulted, 2017. \href{https://arxiv.org/abs/1710.07009}{arXiv:1710.07009}.
\bibitem{lecam} L. Le Cam, \emph{Asymptotic Methods in Statistical Decision Theory}, Springer, New York, 1986. \href{https://doi.org/10.1007/978-1-4612-4946-7}{doi:10.1007/978-1-4612-4946-7}.
\bibitem{torgersen} E. Torgersen, \emph{Comparison of Statistical Experiments}, Cambridge University Press, 1991. \href{https://doi.org/10.1017/CBO9780511666353}{doi:10.1017/CBO9780511666353}.
\bibitem{fritz} T. Fritz, \emph{A synthetic approach to Markov kernels, conditional independence and theorems on sufficient statistics}, Adv. Math. \textbf{370} (2020), 107239. \href{https://doi.org/10.1016/j.aim.2020.107239}{doi:10.1016/j.aim.2020.107239}.
\bibitem{zhu-moran} S. Zhu, \emph{Asymptotic uniformity of the quantization error for Moran measures on $\mathbb R^1$}, preprint, 2018. \href{https://arxiv.org/abs/1802.03723}{arXiv:1802.03723}.
\bibitem{nonanticipative} C. D. Charalambous, P. A. Stavrou, and N. U. Ahmed, \emph{Nonanticipative rate distortion function and relations to filtering theory}, IEEE Trans. Automat. Control \textbf{59} (2014), 937--952. \href{https://doi.org/10.1109/TAC.2013.2293403}{doi:10.1109/TAC.2013.2293403}.
\bibitem{mcdonald} C. McDonald and S. Y\"uksel, \emph{Exponential filter stability via Dobrushin\textquotesingle s coefficient}, arXiv:1910.08463v2 (2020). \href{https://arxiv.org/abs/1910.08463v2}{Version consulted}.
\bibitem{kara} A. D. Kara and S. Y\"uksel, \emph{Near optimality of finite memory feedback policies in partially observed Markov decision processes}, J. Mach. Learn. Res. \textbf{23} (2022), no. 11, 1--46. \href{https://jmlr.org/papers/v23/20-1152.html}{Publisher record}.
\bibitem{supplement} Q. Qi, \emph{General Theta Foundations I: Acquired Geometry and Causal Resource Transfer}, companion research edition, 2026; accompanying mathematical supplement.
\end{thebibliography}

\end{document}
